\chapter{HIV Viral Load}

\section*{What Is This Test?}
HIV viral load (HIV RNA quantitation) measures the amount of HIV-1 RNA in plasma, expressed as copies per milliliter (copies/mL). It is the principal laboratory test for monitoring response to antiretroviral therapy (ART), confirming treatment failure, detecting virologic rebound, and assessing the risk of HIV transmission. Viral load is also used as a diagnostic test for acute HIV infection (HIV-1 RNA positive in the absence of HIV antibodies). Currently available FDA-approved assays use real-time reverse-transcription polymerase chain reaction (RT-PCR) on platforms including Roche cobas 6800/8800 (limit of detection/LOD 20 copies/mL), Abbott RealTime HIV-1 (LOD 40 copies/mL), and Hologic Aptima HIV-1 Quant Dx (LOD 30 copies/mL). The goal of ART is to achieve and maintain virologic suppression below the lower limit of detection (typically <20--50 copies/mL depending on assay). Virologic failure is defined as confirmed HIV RNA >200 copies/mL in a patient on ART for >6 months or failure to achieve suppression by 24 weeks. ART reduces HIV RNA to undetectable levels in the vast majority of adherent patients. The landmark HPTN 052, PARTNER, and Opposites Attract studies have conclusively demonstrated that individuals with sustained undetectable viral loads (<200 copies/mL) do not sexually transmit HIV---the U=U (Undetectable = Untransmittable) principle endorsed by the CDC since 2017 \cite{rodger2019risk}.

\section*{Alternative Names}
HIV RNA Quantitative PCR; HIV Viral Burden; HIV-1 Viral Load; Plasma HIV RNA; HIV RNA Quantitation; HIV PCR; HIV NAT; HIV Nucleic Acid Amplification Test

\section*{Principle}
Real-time reverse-transcription polymerase chain reaction (RT-PCR) amplifies conserved regions of the HIV-1 genome (gag and LTR regions for Roche cobas; pol integrase for Abbott RealTime; 5' LTR and pol for Hologic Aptima). The sample is spiked with a known quantity of an internal quantitation standard (QS, a non-HIV RNA construct with primer-binding sites identical to the HIV target) that controls for extraction and amplification efficiency. HIV-1 RNA copies are calculated by comparing the HIV signal to the QS signal. The assays detect HIV-1 group M subtypes A--H, group O, and group N (Roche) with equivalent quantitation across subtypes. Results reported as copies/mL (converted to log$_{10}$ copies/mL for clinical use: 1 log$_{10}$ change represents a 10-fold change). The lower limit of quantitation/LLOQ ranges from 20 to 40 copies/mL \cite{sollis2014systematic}. Results between the limit of detection and LLOQ are reported as "HIV-1 RNA detected, <LLOQ." The dynamic range extends to approximately 10$^7$ copies/mL. Per the CDC algorithm, qualitative HIV-1 RNA assays (APTIMA HIV-1 RNA Qualitative Assay) are used for diagnostic confirmation, while quantitative assays are used for monitoring \cite{eshleman2014analysis}.

\section*{History}
HIV-1 RNA quantitation first became clinically feasible in the mid-1990s with the development of branched DNA (bDNA) signal amplification (Quantiplex, Chiron) and RT-PCR (Amplicor, Roche). The landmark ACTG 175 and Delta trials (1995--1996) demonstrated that changes in HIV-1 RNA predicted clinical progression, establishing viral load as the key surrogate marker for treatment efficacy. The goal of "undetectable" viral load emerged with more sensitive assays: first-generation assays detected >400 copies/mL; second-generation assays (ultrasensitive RT-PCR) achieved detection limits of 50 copies/mL by the late 1990s; current 4th-generation assays detect 20 copies/mL. The 2014 DHHS guidelines adopted HIV-1 RNA testing for both diagnostic confirmation and treatment monitoring \cite{dhhs2023guidelines}. The U=U campaign, launched in 2016 by the Prevention Access Campaign and endorsed by the CDC in 2017, fundamentally changed the public health messaging and the role of viral load beyond individual treatment to include prevention.

\section*{Why This Test Is Ordered}
\begin{itemize}
  \item Baseline HIV-1 RNA quantitation at the time of HIV diagnosis (entry into care), before ART initiation: establishes the pre-treatment viral load set point, which is prognostic (higher set points are associated with more rapid disease progression)
  \item Monitoring of the virologic response to ART: HIV-1 RNA measured at 4--8 weeks after ART initiation, then every 4--8 weeks until virologic suppression is achieved (<20--50 copies/mL), then every 3--6 months for patients with durable suppression
  \item Detection of virologic failure: defined as confirmed HIV-1 RNA >200 copies/mL after achieving virologic suppression, or failure to achieve viral load below the LLOQ within 24 weeks of ART initiation
  \item Assessment prior to simplification of ART (e.g., switching from 3-drug to 2-drug regimen) or consideration of long-acting injectable ART (cabotegravir/rilpivirine)
  \item Evaluation of acute (primary) HIV infection: HIV-1 RNA is the most sensitive test for acute HIV infection (detectable within 7--10 days of infection), before p24 antigen or antibodies appear; used in the CDC diagnostic algorithm when the differentiation assay is indeterminate
  \item Monitoring of HIV-positive pregnant women at the time of labor and delivery: viral load near delivery determines the mode of delivery (cesarean section recommended if HIV-1 RNA >1,000 copies/mL near delivery) and the infant's ARV prophylaxis regimen
  \item Assessment for participation in research protocols involving analytical treatment interruptions (ATI)
\end{itemize}

\section*{Primary vs Secondary Test}
HIV viral load is the primary test for monitoring response to ART and is one component of the diagnostic algorithm for acute HIV infection. It is NOT used as a screening test (the 4th-generation antigen/antibody assay is the screening test). Once HIV diagnosis is confirmed, the pre-treatment HIV-1 RNA and CD4 count together constitute the essential baseline labs.

\section*{How This Test Is Performed}
Venous blood is collected in an EDTA (purple-top) or plasma preparation tube (PPT, white-top). Plasma must be separated from cells within 6 hours of collection (if stored at room temperature) or within 24 hours (if refrigerated at 2--8\textdegree{}C). Plasma is frozen at -20\textdegree{}C to -80\textdegree{}C if testing is delayed. Separated plasma is stable at -20\textdegree{}C for 5 days or -70\textdegree{}C for extended periods. For the cobas 6800/8800, approximately 1 mL of plasma is required. The presence of heparin in the sample (green-top tube) inhibits PCR and must be avoided. Standard universal precautions for venipuncture apply.

\section*{What Preparation Is Needed}
No patient fasting or special preparation required. Timing of the blood draw relative to the last ART dose is not critical for most assays; however, consistent timing is recommended for intra-patient comparisons.

\section*{Contraindications and Precautions}
HIV-1 viral load assays may not reliably quantitate HIV-2 RNA, as these assays are specifically designed for HIV-1. If an HIV-2-infected patient requires viral load monitoring, an HIV-2-specific assay must be arranged through a reference laboratory or the CDC. HIV-1 RNA results may be falsely elevated in the setting of intercurrent illness (e.g., acute infection with another pathogen, recent vaccination), which can cause transient viremia ("blips"---isolated HIV RNA between the LLOQ and 500 copies/mL that returns to undetectable on repeat testing without any change in ART). A confirmed viral load >200 copies/mL (not just a "blip") is the threshold for virologic failure. "Blips" (isolated low-level viremia 20--500 copies/mL) are not predictive of virologic failure and do not require a change in ART. Hemolyzed specimens, lipemic specimens, and specimens collected in heparin tubes may interfere with amplification. Plasma must be separated from cells promptly to prevent HIV RNA degradation by cellular nucleases.

\section*{Risks}
Standard venipuncture risks: pain, bruising, hematoma, phlebotomy site infection, vasovagal syncope. No risks from the assay itself.

\section*{What Happens After the Test}
HIV-1 RNA results are typically available in 24--72 hours. Results guide ART management: if the viral load is declining appropriately from baseline (targeting undetectable by 12--24 weeks), the current ART regimen is continued with monitoring every 3--6 months. If virologic failure is confirmed (HIV-1 RNA >200 copies/mL on two consecutive measurements >2 weeks apart), resistance testing should be performed while the patient is on the failing regimen (or within 4 weeks of discontinuation) \cite{saag2020antiretroviral}. The ART regimen should then be changed based on resistance testing results, adherence assessment, and tolerability. For patients with sustained virologic suppression and stable clinical status, viral load monitoring can be extended to every 6 months.

\section*{Understanding Results}

\subsection*{Below Normal}
\begin{itemize}
  \item \textbf{HIV-1 RNA not detected (target not detected, TND):} Viral RNA is below the assay's limit of detection (<20 copies/mL for cobas). This is the goal of ART. Indicates excellent adherence and potent regimen. Associated with zero risk of sexual transmission of HIV (U=U). Used to qualify for simplification strategies, including dual-drug ART and long-acting injectable therapy
  \item \textbf{HIV-1 RNA detected, <LLOQ (<20 or <40 copies/mL):} Extremely low-level viremia at the limit of detection. Not considered virologic failure. Reassuring; continue current ART and monitoring
  \item \textbf{Negative HIV-1 RNA in a patient not on ART:} In the diagnostic algorithm, a reactive fourth-generation screen followed by a negative or indeterminate differentiation assay and negative HIV-1 NAT generally indicates no HIV-1 infection. If acute infection, HIV-2, recent exposure, or discordant clinical findings are possible, follow the complete current diagnostic algorithm and arrange repeat or supplemental testing.
\end{itemize}

\subsection*{Above Normal}
\begin{itemize}
  \item \textbf{HIV-1 RNA detectable at >20 (or >40) but <200 copies/mL:} Low-level viremia. If isolated ("blip"), no intervention required; repeat in 2--4 weeks. If persistent on multiple draws (>2 consecutive values), evaluate adherence, drug-drug interactions, absorption issues, and consider resistance testing if >200 copies/mL
  \item \textbf{HIV-1 RNA >200 copies/mL (confirmed on two measurements):} Virologic failure. Resistance testing should be performed, and ART should be changed based on genotype results, treatment history, and tolerability. This threshold (>200) is used to define virologic failure in clinical trials and guidelines. Prompt evaluation and intervention reduce the risk of accumulating resistance mutations
  \item \textbf{HIV-1 RNA >100,000 copies/mL at baseline (pre-ART):} High viral set point associated with more rapid CD4 decline and higher risk of HIV-associated complications. No specific ART choice is mandated, but INSTI-based regimens (bictegravir/FTC/TAF, dolutegravir/ABC/3TC or DTG + FTC/TAF) are preferred first-line due to rapid virologic suppression and high genetic barrier to resistance
  \item \textbf{HIV-1 RNA >1,000 copies/mL in late pregnancy:} Indicates high risk of perinatal transmission (MTCT rate increases significantly). Cesarean delivery at 38 weeks (before onset of labor and rupture of membranes) is recommended. The neonate should receive combination ARV prophylaxis presumptively
\end{itemize}

\section*{Normal Results}
For an HIV-uninfected individual: \textbf{HIV-1 RNA not detected (target not detected)}. For an HIV-infected individual on ART with virologic suppression: \textbf{HIV-1 RNA not detected (target not detected)} or \textbf{<20 copies/mL} or \textbf{<40 copies/mL} depending on assay. Note that "normal" means suppressed, not absent, because HIV persists as integrated proviral DNA in latently infected resting CD4+ T cells (the viral reservoir) despite ART.

\section*{Follow-up Tests}
\begin{itemize}
  \item \textbf{Repeat HIV-1 RNA in 2--4 weeks:} For isolated low-level viremia (20--200 copies/mL) to distinguish a "blip" from early virologic failure. Two consecutive values >200 copies/mL define failure
  \item \textbf{HIV drug resistance genotype (protease, reverse transcriptase, integrase):} Should be obtained while the patient is on the failing ART regimen or within 4 weeks of discontinuation. Guides selection of the next ART regimen. Accumulated resistance mutations from a failing NNRTI-based regimen (particularly efavirenz or rilpivirine) can limit future options
  \item \textbf{CD4 T-lymphocyte count:} Performed concurrently with viral load at baseline and for monitoring; CD4 count guides OI prophylaxis and immune reconstitution assessment
  \item \textbf{HIV-2 viral load:} For HIV-2-infected patients; an HIV-2-specific quantitative RNA or DNA assay performed at a reference laboratory (University of Washington, CDC)
  \item \textbf{ART therapeutic drug monitoring (TDM):} Not routinely recommended but may be considered in select situations (suspected malabsorption, drug-drug interactions, pregnancy physiology altering pharmacokinetics, or failure of a protease inhibitor- or NNRTI-based regimen)
\end{itemize}

\section*{References}
\bibliographystyle{unsrtnat}
\bibliography{references}

\clearpage
\section*{Regulatory Status and Authorization Date}
Regulatory status applies to the specific assay, instrument, manufacturer, and intended use rather than to the test name alone. No single approval date can be assigned to this test category. For a commercial assay, verify the current product in the relevant regulator's database: the U.S. Food and Drug Administration (FDA) clearance, approval, or authorization; India's Central Drugs Standard Control Organisation (CDSCO) licence or permission; the European Union In Vitro Diagnostic Regulation (EU IVDR) CE certificate issued through the applicable conformity-assessment route; or the United Kingdom Medicines and Healthcare products Regulatory Agency (MHRA) registration and UKCA/CE status. Laboratory-developed tests may not have an individual FDA, CDSCO, EU IVDR, or MHRA approval date.
\clearpage
